\section{Isometries, Unitary Operators, and Matrix Factorization}
\subsection{Isometries}

\begin{mydef}[isometry]
  $S \in \lin{V}{W}$ is called an \qt{isometry} if
  \[
  \norm{Sv} = \norm{v} \quad \forall v \in V.
  \]

  In other words, a linear map is an isometry if it preserves norms.

  If $S \in \linmap(V,W)$ is an isometry and $v \in V$ is \st $Sv = 0$, then
  \[
  \norm{v} = \norm{Sv} = \norm{0} = 0,
  \]

  which implies that $v=0$. Thus, every isometry is injective.
\end{mydef}

(The Greek word \emph{isos} means \qt{equal}; the Greek word \emph{metron} means \qt{measure}. Thus \emph{isometry} literally means \qt{equal measure}.)

\mce{49}
\begin{thm}[characterizations of isometries]
  \label{thm: characterizations of isometries}
  Suppose $S \in \linmap (V,W)$. Suppose $e_1, \ddd, e_n$ is an orthonormal basis of $V$ and $f_1, \ddd, f_m$ is an orthonormal basis of $W$. Then the following are equivalent:
  \begin{enumerate}[label=\textbf{(\alph*)}]
    \item $S$ is an isometry;
    \item $S^*S = I$;
    \item $\ip{Su}{Sv} = \ip{u}{v} \quad \forall u,v \in V$;
    \item $Se_1, \ddd, Se_n$ is an orthonormal list in $W$;
    \item The columns of $\mmatrix(S, (e_1, \ddd, e_n), (f_1, \ddd, f_m))$ form an orthonormal list in $\myF^m$ with respect to the Euclidean inner product.
  \end{enumerate}
\end{thm}
\begin{prf}
  We will prove $(a) \implies (b) \implies (c) \implies (d) \implies (e) \implies (a)$.

  \qt{First step:} First, suppose (a) holds, so $S$ is an isometry. If $v\in V$, then:
  \[
  \ip{(I-S^*S)v}{v} = \ip{v}{v} - \ip{S^*Sv}{v} = \norm{v}^2-\ip{Sv}{Sv} = \norm{v}^2 - \norm{Sv}^2 = 0.
  \]

  Since $I-S^*S$ is a self-adjoint operator because $(I-S^*S)^* = I-S^*S$, we can use \ref{thm: T self-adjoint and <Tv,v>=0 for all v <=> T=0} to conclude that $I-S^*S$ equals the zero operator. Thus, $S^*S = I$, proving $(a) \implies (b)$.

  \qt{Second step:} Now, suppose (b) holds, so $S^*S=I$. If $u,v \in V$, then:
  \[
  \ip{Su}{Sv} = \ip{S^*Su}{v} = \ip{Iu}{v} = \ip{u}{v},
  \]
  proving $(b) \implies (c)$.

  \qt{Third step:} Now, suppose that (c) holds, so $\ip{Su}{Sv} = \ip{u}{v} \quad \forall u,v \in V$. Thus, if $j,k \in \{1, \ddd, n\}$, then:
  \[
  \ip{Se_j}{Se_k} = \ip{e_j}{e_k} = \begin{cases} 1, & \text{if } j = k, \\ 0, & \text{if } j \neq k. \end{cases}
  \]

  Hence, $Se_1, \ddd, Se_n$ is an orthonormal list in $W$, proving that $(c) \implies (d)$.

  \qt{Fourth step:} Suppose (d) holds. The $k$-th column of $\mmatrix(S) \equiv \mmatrix(S, (e_1, \ddd, e_n), (f_1, \ddd, f_m))$ consists of the scalars $(\ip{Se_k}{f_1}, \ddd, \ip{Se_k}{f_m})^t$. By \ref{thm: writing a vector as a linear combination of an orthonormal basis}(c), the Euclidean inner product in $\myF^m$ of the $j$-th and $k$-th columns of $\mmatrix(S)$ equals
  \[
    \sum_{\ell=1}^m \ip{Se_j}{f_\ell} \overline{\ip{Se_k}{f_\ell}} = \ip{Se_j}{Se_k} = \begin{cases} 1, & \text{if } j = k, \\ 0, & \text{if } j \neq k. \end{cases}
  \]
  Thus, the columns of $\mmatrix(S)$ form an orthonormal list in $\myF^m$, proving $(d) \implies (e)$.

  \qt{Fifth step:} Suppose (e) holds. Then the vectors $Se_1, \ddd, Se_n$ form an orthonormal list in $W$. For any $v \in V$, write $v = a_1 e_1 + \cdots + a_n e_n$ for some $a_1, \ddd, a_n \in \myF$. Then $Sv = a_1 Se_1 + \cdots + a_n Se_n$. By \ref{thm: norm of an orthonormal linear combination}:
  \[
    \norm{Sv}^2 = |a_1|^2 + \cdots + |a_n|^2 = \norm{v}^2.
  \]
  Taking square roots gives $\norm{Sv} = \norm{v}$ for all $v \in V$, so $S$ is an isometry, proving $(e) \implies (a)$.
\end{prf}

\subsection{Unitary Operators}

\mce{51}
\begin{mydef}[unitary operator]
  An operator $S \in \linmap (V)$ is called \qt{unitary} if $S$ is an invertible isometry.
\end{mydef}

Note: Every isometry is injective. Every injective operator on a finite-dimensional vector space is invertible. So \qt{unitary} and \qt{isometry} mean the same thing for operators on finite-dimensional inner product spaces. But remember: A unitary operator maps a vector space to itself, while an isometry maps a vector space to another (possibly different) vector space.

\begin{example}[rotation of $\real^2$]
  Suppose $\theta \in \real$ and $S$ is the operator on $\myF^2$ whose matrix $\mmatrix(S)$ with respect to the standard basis of $\myF^2$ is
  \[
    \left(\begin{array}{cc}
      \cos \theta & - \sin \theta \\
      \sin \theta & \cos \theta
    \end{array} \right)
    = \mmatrix(S).
  \]

  The two columns of this matrix form an orthonormal list in $\myF^2$. Therefore, $S$ is an isometry by the equivalence of (a) and (e) in \ref{thm: characterizations of isometries}. Thus, $S$ is a unitary operator.

  If $\myF=\real$, then $S$ is the counterclockwise rotation by $\theta$ radians around the origin of $\real^2$. So it preserves norms.
\end{example}

\mce{53}
\begin{thm}[characterizations of unitary operators]
  \label{thm: characterizations of unitary operators}
  Suppose $S \in \linmap(V)$ and $e_1, \ddd, e_n$ is an orthonormal basis of $V$. Then the following are equivalent:
  \begin{enumerate}[label=\textbf{(\alph*)}]
    \item $S$ is a unitary operator;
    \item $S^*S = SS^* = I$;
    \item $S$ is invertible and $S^{-1} = S^*$;
    \item $Se_1, \ddd, Se_n$ is an orthonormal basis of $V$;
    \item The rows of $\mmatrix(S,(e_1, \ddd, e_n))$ form an orthonormal basis of $\myF^n$ with respect to the Euclidean inner product;
    \item $S^*$ is a unitary operator.
  \end{enumerate}
\end{thm}
\begin{prf}
  An operator $S \in \linmap(V)$ on a finite-dimensional space is an isometry $\iff S^*S = I$ (by \ref{thm: characterizations of isometries}).
  Since $V$ is finite-dimensional, $S^*S = I \iff SS^* = I \iff S^{-1} = S^*$, which establishes the equivalence of (a), (b), and (c).
  The equivalence with (d) follows because $S$ is an isometry $\iff Se_1, \dots, Se_n$ is an orthonormal list of length $\dim V$, hence an orthonormal basis.
  The equivalence with (e) and (f) follows because $S^*$ is unitary $\iff (S^*)^* S^* = S S^* = I$, and the rows of $\mmatrix(S)$ are the conjugates of the columns of $\mmatrix(S^*)$.
\end{prf}

\mce{54}
\begin{thm}[eigenvalues of unitary operators]
  \label{thm: eigenvalues of unitary operators}
  Suppose $\lambda$ is an eigenvalue of a unitary operator $S \in \linmap(V)$. Then $|\lambda| = 1$.
\end{thm}
\begin{prf}
  Let $v \neq 0$ be an eigenvector of $S$ corresponding to $\lambda$, so $Sv = \lambda v$.
  Because $S$ is an isometry, $\norm{v} = \norm{Sv} = \norm{\lambda v} = |\lambda| \norm{v}$.
  Since $\norm{v} \neq 0$, dividing by $\norm{v}$ gives $|\lambda| = 1$.
\end{prf}

\mce{55}
\begin{thm}[description of unitary operators on complex inner product spaces]
  \label{thm: description of unitary operators on complex inner product spaces}
  Suppose $\myF = \compl$ and $S \in \linmap(V)$. Then the following are equivalent:
  \begin{enumerate}[label=\textbf{(\alph*)}]
    \item $S$ is a unitary operator;
    \item There is an orthonormal basis of $V$ consisting of eigenvectors of $S$ whose corresponding eigenvalues all have absolute value $1$.
  \end{enumerate}
\end{thm}
\begin{prf}
  \Rightarrowdirection Suppose $S$ is unitary. Then $S^*S = I = SS^*$, so $S$ is normal. By the Complex Spectral Theorem (\ref{thm: complex spectral theorem}), $V$ has an orthonormal basis of eigenvectors of $S$. By \ref{thm: eigenvalues of unitary operators}, each corresponding eigenvalue has absolute value $1$.

  \Leftarrowdirection If $e_1, \dots, e_n$ is an orthonormal basis of eigenvectors with $Se_k = \lambda_k e_k$ and $|\lambda_k| = 1$, then for any $v = \sum a_k e_k$, $\norm{Sv}^2 = \sum |\lambda_k a_k|^2 = \sum |a_k|^2 = \norm{v}^2$. Thus $S$ is an isometry, hence a unitary operator.
\end{prf}

\mce{56}
\subsection{QR Factorization}
\begin{mydef}[unitary matrix]
  An $n$-by-$n$ matrix is called \qt{unitary} if its columns form an orthonormal list in $\myF^n$.
\end{mydef}

\begin{thm}[characterizations of unitary matrices]
  Let $Q \in \myF^{n,n}$. Then the following are equivalent:
  \begin{enumerate}[label=\textbf{(\alph*)}]
    \item $Q$ is a unitary matrix;
    \item The rows of $Q$ form an orthonormal list in $\myF^{n}$;
    \item $\norm{Qv} = \norm{v} \quad \forall v \in \myF^n$;
    \item $Q^*Q=QQ^*=I$, the $n$-by-$n$ matrix with $1$'s on the diagonal and $0$'s elsewhere.
  \end{enumerate}
\end{thm}
\begin{prf}
  This is the matrix translation of \ref{thm: characterizations of unitary operators} applied to the operator $x \mapsto Qx$ on $\myF^n$.
\end{prf}

\begin{thm}[QR factorization]
  \label{thm: QR factorization}
  Suppose $A \in \myF^{n,n}$ is an invertible square matrix. Then there exist unique matrices $Q$ and $R$ such that $Q$ is unitary, $R$ is upper triangular with only positive numbers on its diagonal, and
  \begin{equation}
    A = QR.
  \end{equation}
\end{thm}
\begin{prf}
  \prooffont{Existence:} Let $a_1, \ddd, a_n$ denote the columns of $A$. Apply the Gram-Schmidt Procedure (\ref{thm: gram-schmidt procedure}) to $a_1, \ddd, a_n$ to obtain an orthonormal basis $q_1, \ddd, q_n$ of $\myF^n$.
  Then each $a_k = R_{1,k} q_1 + \cdots + R_{k,k} q_k$ with $R_{k,k} = \norm{a_k - \sum_{j=1}^{k-1} \ip{a_k}{q_j}q_j} > 0$.
  Setting $Q$ to be the unitary matrix with columns $q_1, \ddd, q_n$ and $R = (R_{j,k})$ gives $A = QR$.

  \prooffont{Uniqueness:} If $A = Q_1 R_1 = Q_2 R_2$, then $Q_2^* Q_1 = R_2 R_1^{-1}$. The matrix $Q_2^* Q_1$ is unitary, while $R_2 R_1^{-1}$ is upper triangular with positive diagonal entries. The only such matrix is $I$, so $R_1 = R_2$ and $Q_1 = Q_2$.
\end{prf}

\subsection{Cholesky Factorization}

\mce{61}
\begin{thm}[positive invertible operator]
  \label{thm: positive invertible operator}
  A self-adjoint operator $T\in \linmap(V)$ is a positive invertible operator $\iff \ip{Tv}{v} > 0 \quad \forall v \in V, v \neq 0$.
\end{thm}
\begin{prf}
  \Rightarrowdirection If $T$ is positive and invertible, $\ip{Tv}{v} \ge 0$. If $\ip{Tv}{v} = 0$, by \ref{thm: T positive and <Tv,v> = 0 => Tv = 0}, $Tv = 0 \implies v = 0$. Thus $\ip{Tv}{v} > 0$ for all $v \neq 0$.

  \Leftarrowdirection If $\ip{Tv}{v} > 0$ for all $v \neq 0$, then $T$ is positive. If $Tv = 0$, then $\ip{Tv}{v} = 0 \implies v = 0$, so $\mynull T = \{0\}$, which implies $T$ is invertible.
\end{prf}

\begin{mydef}[positive definite]
  A matrix $B \in \myF^{n,n}$ is called \qt{positive definite} if $B^* = B$ and
  \begin{equation}
    \ip{Bx}{x} > 0 \quad \forall x \in \myF^n, \, x \neq 0.
  \end{equation}
\end{mydef}

\begin{thm}[Cholesky factorization]
  \label{thm: cholesky factorization}
  Suppose $B \in \myF^{n,n}$ is a positive definite matrix. Then there exists a unique upper-triangular matrix $R$ with only positive numbers on its diagonal \st
  \begin{equation}
    B = R^*R.
  \end{equation}
\end{thm}
\begin{prf}
  \prooffont{Existence:} By \ref{thm: uniqueness of square root}, the positive definite matrix $B$ has a unique positive square root $\sqrt{B}$, which is invertible.
  Applying QR factorization (\ref{thm: QR factorization}) to $\sqrt{B}$ gives $\sqrt{B} = QR$ for unitary $Q$ and upper-triangular $R$ with positive diagonal entries.
  Then $B = (\sqrt{B})^* \sqrt{B} = (QR)^* (QR) = R^* Q^* Q R = R^* R$.

  \prooffont{Uniqueness:} If $B = R_1^* R_1 = R_2^* R_2$, then $(R_2 R_1^{-1})^* (R_2 R_1^{-1}) = I$, so $R_2 R_1^{-1}$ is unitary and upper triangular with positive diagonal entries. Hence $R_2 R_1^{-1} = I \implies R_1 = R_2$.
\end{prf}

